\begin{abstract}
Surrogate-assisted optimization is becoming an important part of energy system modeling as high temporal resolution, sector coupling, uncertainty, and multiple objectives increase computational requirements. Yet the literature is split across model families, energy applications, and optimization communities. Existing reviews generally address only parts of this landscape and do not provide a common structure for comparing how surrogates are trained, integrated, and validated. This review develops such a structure for surrogate-assisted energy system optimization, with particular attention to multi-energy systems and multi-objective problems. A systematic search of 3,129 records, an in-depth analysis of 285 studies, a bibliometric analysis, and a meta-review of 32 previous reviews are combined in one framework. The main contribution is a new taxonomy that follows the complete surrogate-assisted optimization workflow: the model class and approximated target, the origin and selection of training data, the role of the surrogate within the optimizer, and the validation of predictions and resulting decisions. It brings together the main surrogate model classes, design-of-experiments and multi-fidelity strategies, five recurring integration patterns, and five application-target domains. These dimensions are not only described but quantified and connected through a machine-readable evidence map, allowing relationships between methods, applications, and validation practices to be examined across the literature. This combined view shows that model choice alone does not explain how surrogates support energy system optimization; their value depends equally on training-data design, integration role, and the evidence used to establish decision reliability. The review therefore provides both a consolidated account of current practice and a basis for method selection, comparison, and future research. It concludes with the software landscape and a research agenda focused on decision-aware benchmarking, calibrated uncertainty, out-of-distribution testing, reusable datasets, uncertainty-aware multi-criteria decision making, and interoperable workflows.
\end{abstract}
